\chapter{ISR spin-correlation corrections}
\label{app:spin-correl}

This appendix derives the additional NNLL correction that arises when
a gluon enters the hard scattering. In addition to the ISR
hard-collinear corrections
$C_{{\rm hc},\ell}^{(1)}$, $\delta\mathcal{F}_{\rm hc}$, and
$\delta\mathcal{F}_{\rm rec}$ derived in
Sec.~\ref{sec:NNLL-initial}, an incoming gluon introduces a further
contribution with no counterpart for quark-initiated radiation: spin
correlations. These give rise to the additional NNLL contribution
$\Delta\mathcal{F}_{\rm rec}$, which vanishes for one-jettiness but is
required for the general ISR formalism and for observables with
non-trivial azimuthal dependence.

To account for this, we introduce the spin-density matrix for the hard process
in $4-2\varepsilon$ dimensions. Neglecting terms that are antisymmetric in $\mu,\nu$, the general form is 
\begin{equation}
\rho^{\mu\nu} = \frac{d^{\mu\nu} }{2(1-\varepsilon)}
+ \mathcal{T}_\ell(\{p\};\varepsilon)\left(2n_{\rm in}^\mu n_{\rm in }^\nu - \frac{d^{\mu\nu} }{1-\varepsilon}\right) \,,
\end{equation}
where we choose $d^{\mu\nu}$ to be the sum over physical polarisations transverse to both $p_1$ and $p_2$,
\begin{equation}
d^{\mu\nu} = -g^{\mu\nu} + \frac{p_1^\mu p_2^\nu + p_2^\mu p_1^\nu}{p_1 \cdot p_2}\,,
\end{equation}
and $n_{\rm in}$ is a space-like unit vector that lies in the plane identified
by the beam and the $Z$ boson, satisfying $n_{\rm in}^2=-1$ and $n_{\rm in}\cdot p_1=n_{\rm in}\cdot p_2=0$. We obtain the following relations: 
\begin{subequations}
  \begin{align}
  &d^{\mu\nu}d_{\mu\nu}=-g^{\mu\nu}d_{\mu\nu}=2(1-\varepsilon)\,,\\[1.5ex]
  &d_{\mu\nu} \left(2n_{\rm in}^\mu n_{\rm in }^\nu - \frac{d^{\mu\nu}}{1-\varepsilon}\right)=-g^{\mu\nu}\left(2 n_{\rm in}^\mu n_{\rm in }^\nu - \frac{d^{\mu\nu}}{1-\varepsilon}\right)=0\,,\\[1.5ex] 
  &\left(2 n_{\rm in}^\mu n_{\rm in }^\nu - \frac{d^{\mu\nu}}{1-\varepsilon}\right)\left(2 n_{{\rm in},\mu} n_{{\rm in },\nu} - \frac{d_{\mu\nu}}{1-\varepsilon}\right) = 2 \frac{1-2\varepsilon}{1-\varepsilon}\,.
  \end{align}
\end{subequations}
This implies
\begin{equation}
  \label{eq:PL}
  \mathcal{T}_\ell(\{\tilde p\};\varepsilon) = \frac{1-\varepsilon}{2(1-2\varepsilon)} \left(2 n_{\rm in}^\mu n_{\rm in }^\nu - \frac{d^{\mu\nu}}{1-\varepsilon}\right) \rho_{\mu\nu}\,.
\end{equation}

The collinear splitting probability is then obtained by contracting
$\rho^{\mu\nu}$ with the unaveraged splitting functions
\begin{equation}
\left[\hat{P}_{ga}^{(1)}(z,\vec k_t;\varepsilon)\right]^{\mu\nu}
= -g^{\mu\nu}\,\hat{P}_{ga}(z;\varepsilon)
+ \left(\frac{g^{\mu\nu}}{1-\varepsilon} + 2\frac{k_\perp^\mu k_\perp^\nu}{k_t^2}\right)
2\hat G_{ga}(z;\varepsilon)\,, \qquad a = g,q\,,
\end{equation}
with $k_\perp\cdot k_\perp = -k_t^2$, and
\begin{equation}
  \label{eq:Grho}
  d_{\mu\nu}\left(\frac{g^{\mu\nu}}{1-\varepsilon} + 2\frac{k_\perp^\mu k_\perp^\nu}{k_t^2}\right)=0\,.
\end{equation}
The azimuthally averaged splitting functions $\hat{P}_{ga}(z;\varepsilon)$ with $a=q,g$ are those present in Eqs.~\eqref{eq:Pgq-epsilon} and.~\eqref{eq:Pgg-epsilon} respectively, and the
spin-correlated functions are 
\begin{equation}
\begin{aligned}
\hat G_{gg}(z) &= -(1-\varepsilon)C_A \frac{1-z}{z}\,,\qquad
\hat G_{gq}(z) &= C_F \frac{1-z}{z}\,.
\end{aligned}
\end{equation}
Performing the contraction, using the fact $k_\perp \cdot n_{\rm in }=-k_t \cos\phi$ we obtain
\begin{equation}
\rho^{\mu\nu}\left[\hat{P}_{ga}(z,\vec k_t;\varepsilon)\right]_{\mu\nu}
= \hat{P}_{ga}(z;\varepsilon)
+  \mathcal{T}_\ell(\{\tilde p\};\varepsilon)\,\hat G_{ga}(z;\varepsilon)\left[2(1-\varepsilon)\cos^2\phi - 1\right]\,,
\qquad a = g,q\,.
\end{equation}
The above equation gives the probability for a splitting $a\to g
a$. The contribution proportional to the term
splitting function $\hat{P}_{ga}(z;\varepsilon)$ can be treated exactly
as in Sec.~\ref{sec:NNLL-initial}, leading to the ISR
hard-collinear coefficient $C_{\mathrm{hc},\ell}^{(1)}$ together with
the corrections $\delta\mathcal{F}_{\rm hc}$ and
$\delta\mathcal{F}_{\rm rec}$. The remaining term, proportional to
$\hat{G}_{ga}(z;\varepsilon)$, contains the spin-correlation
dependence. Following the same derivation as
Ref.~\cite{Arpino:2019ozn}, this gives rise to the additional NNLL
correction
\ABquestion{Is this correct sum over legs here? Do all the $\ell$'s need to be $\ell_g$
and does the $z$ and $\phi$ need to be $z^{\ell_g}$ and $\phi^{\ell_g}$}
\begin{equation}
\begin{aligned}
\label{eq:DFrec-ISR}
  \Delta&\mathcal{F}_{\mathrm{rec}}(\lambda) =
\sum_{\ell_g}\frac{\alpha_s(Qv^{1/(a+b_{\ell_g})})}{\alpha_s(Q)(a+b_{\ell_g})}\int_0^\infty \frac{d\zeta}{\zeta}
\int_0^\pi \frac{d\phi}{\pi}\left[2\cos^2\phi^{(\ell_g)} - 1\right]
\int_0^1\!dz\,
\Delta P_g(z^{(\ell_g)})
\\
&\quad \times \int \dZ \left[
\Theta\left(1 - \frac{V_{\mathrm{hc}}(\{\tilde{p}\},k,\{k_i\})}{v}\right)
- \Theta(1-\zeta)\,
\Theta\left(1 - \frac{V_{\mathrm{sc}}(\{\tilde{p}\},\{k_i\})}{v}\right)
\right]\,,
\end{aligned}
\end{equation}
where the sum runs over the gluon legs $\ell_g$. For initial-state
radiation,
\begin{equation}
\label{eq:DPg-ISR}
\Delta P_g(z) =
4\left(\frac{1-z}{z}\right)\frac{\Theta(z-x_\ell)}{z f_g(x_\ell,Q v^{1/(a+b_\ell)})}
\left[C_A f_g\left(\frac{x_\ell}z,Q v^{1/(a+b_\ell)}\right) + C_F f_q\left(\frac{x_\ell}z,Q v^{1/(a+b_\ell)}\right)\right]\,,
\end{equation}
while for final-state radiation
\begin{equation}
\label{eq:DPg-FSR}
\Delta P_g(z) =4z(1-z)\left[\frac{C_A}{2}-T_Rn_f\right]\,.
\end{equation}